% ============================================================
% 第八章 · 小组总结
% 对应大纲：# 八、小组总结
% ============================================================

\chapter{小组总结}


9.1 项目开发全景复盘：从 0 到 1 的协作落地
在“程序设计与制作”课程大作业的驱动下，我们“六个基米哈一天”小组以《第
五人格》为 IP 原型，耗时三周完成网页小游戏《第五人格：推演终幕》的开发。项目以
“求生者剧情推演 + 监管者卡牌战斗”为核心玩法，涵盖剧情系统、战斗系统、存档系
统、地图交互等六大核心模块，最终交付包含 3 万字剧情、17 个人物档案、2 张可切换
地图、1 个修机小游戏及豆瓣电影 / 中华美食两个附属网站的完整成果。回顾开发全流
程，我们以 “每日会议为协作枢纽、分工互补为推进引擎、问题共解为突破关键”，实
现了从需求拆解到功能落地的高效推进。
9.1.1 第一周：框架搭建与基础攻坚（8.24-8.30）
项目启动阶段，我们聚焦 “定方向、搭框架、破初障”三大目标。8 月 24 日 - 25
日，组长赵硕涵牵头完成组队与方向决策，组织成员学习项目 PPT 后，经线下会议投票
确定开发“《第五人格》IP 二创交互游戏”，明确“剧情 + 战斗”双核心玩法；房金衡作
为剧情与战斗策划，整合成员创意形成第一版游戏大纲，规划“序章 + 四章主线”剧情
结构，定义“求生者推演（剧情向）+ 监管者战斗（卡牌向）”的核心循环；梁思齐、刘
雨桐、潘治鑫组成技术组，同步启动前端基础学习，王宇豪则单美术线推进，搭建像素
素材库。8 月 26 日 - 30 日，各模块进入基础开发阶段。技术组方面，梁思齐从前端三件
套入门，克服 CSS 语法误区（误写 tiles 类后改用 button 标签），完成登录界面与角色
移动类初步开发；刘雨桐聚焦成就系统，参考往届“肖申克”“瑞雪”等作品，拆解出
“Achievements.json 配置表 + AchievementManager 管理逻辑 + 弹窗 UI”的模块架构，
同步编写剧情对话界面；潘治鑫则钻研 localStorage 存储逻辑，为存档系统奠定基础。
剧情与战斗组方面，房金衡完成序章“初入庄园”剧本及第一章前两幕剧情，同时设计
战斗系统基础规则，定义 90 点初始生命值、2 点行动点的玩家参数，规划普通攻击、利
剑之刺等 12 类基础卡牌。美术组方面，王宇豪从零学习 Tiled Map Editor 与 Aseprite，完
成首个测试地图与主角像素小人初稿，解决素材风格不统一问题。此阶段，我们遭遇首
个技术难关——CORS 协议跨域问题。8 月 29 日，梁思齐、潘治鑫在调试地图切换与存
档加载时，发现直接打开 HTML 文件无法获取数据，经每日会议集体排查，最终参考往
届方案，通过 Python HTTP 服务器与 Live Server 插件双路径解决，同时总结出“避免绝
对路径、统一相对路径调用”的开发规范，为后续模块协作扫清障碍。
9.1.2 第二周：功能深化与协同突破（8.31-9.6）
进入核心开发阶段，我们以 “模块协同、问题共解”为核心，推动项目从“能用”
向“好用”升级。技术组实现多系统联动：梁思齐优化角色移动逻辑，修复“对话未完
成关闭后状态异常”“空气墙不精准”等 5 类 Bug，同时创新性地用 Web Audio API 将
地图钢琴改造为可交互小游戏，解决“Click Noise”杂音问题；刘雨桐重构剧情系统架
构，将 20 个场景剧情拆分至独立 JS 文件，通过 scriptEngine.js 统一接口，解决浏览
器加载限制（加入 datamanager 与 UCL 参数加载函数）；潘治鑫则分离成就与存档系
42
项目计划
统（避免前期耦合导致的维护难题），新增自动存档功能，实现“剧情跳转后角色位置
与背包状态无损恢复”。剧情与战斗组实现内容扩容：房金衡完成第一章收尾（厂长剧
情）与第二章审核，同步优化战斗系统，为监管者新增特色机制（如红夫人“重生”、蜘
蛛“蛛网反伤”），调整卡牌数值避免流派失衡；同时拓展附属项目，耗时 10 余小时用
Python（requests+lxml）爬取豆瓣电影 Top250 数据，突破反爬限制（设置 1 秒延迟、伪
造 headers），将电影名称、评分、海报等信息以 JSON 格式存储，完成电影欣赏网站开
发。美术组则推进地图迭代：王宇豪完成地下室地图第二版，优化房间布局（划分图书
区、药水区），通过 GIMP 缩放处理 16×16 与 32×32 素材比例问题，同时绘制带感叹号
的交互贴图，提升场景引导性。此阶段的关键突破在于 “跨模块协作机制”的建立。9
月 2 日，梁思齐在开发钢琴小游戏时，发现 iframe 与角色移动的鼠标聚焦冲突，经与王
宇豪（美工）、潘治鑫（存档）召开专项会议，最终实现“点击游戏区域控制小游戏、点
击下方控制角色”的分层交互逻辑；9 月 4 日，刘雨桐因数学建模竞赛请假前，提前完
成剧情系统文档梳理，确保潘治鑫能顺利对接存档接口，体现“无缝交接”的协作意识。
9.1.3 第三周：整合优化与成果交付（9.7-9.13）
收尾阶段，我们聚焦 “全流程测试、细节优化、成果整合”，确保项目符合课程交
付标准。技术组完成系统闭环：梁思齐解决战斗系统与存档的数据断联问题（通过日志
追踪定位参数传递时机错误），实现“战斗胜利后自动触发剧情、失败后重置状态”的
完整流程；潘治鑫开发修机小游戏，设计“校准进度条 + 钥匙解锁门”的地图衔接逻辑，
同时完成成就系统与结局的联动（Boss 战胜利触发“欧律狄刻”成就，失败触发“故事
的最后”成就）；刘雨桐则完成角色档案系统开发，通过 localStorage 实现“账号级解
锁状态存储”，同步导入全部 3 万字剧情，解决文件命名格式错误（如“1.jpg”误判为二
进制文件）。剧情与美术组完成内容收尾：房金衡撰写第四章双结局（“精神清醒”“永
坠混沌”），补充 17 个人物档案（含生日、特质、背景故事），同时调整战斗数值（如
作曲家“音韵”被动技能改为“回合开始抽 1 张牌”）；王宇豪则对地下室地图进行像
素级优化，为关键道具（实验笔记）添加 Aseprite 发光动画，校准所有碰撞框，确保场
景交互无死角。最终交付阶段，我们建立 “模块测试 - 交叉验收 - 整体联调”的三级质
检体系：赵硕涵统筹全流程测试，重点验证地图切换、存档加载、成就触发等关键路径；
各成员交叉测试非负责模块（如梁思齐测试剧情跳转、王宇豪测试战斗 UI），最终解决
“非 Live Server 环境钥匙获取异常”“结局页面音效不同步”等 8 类遗留问题，确保项目
稳定运行。
9.2 技术难关攻克：从“卡壳”到“破局”的能力跃迁
三周开发中，我们遭遇的技术挑战不仅是“问题清单”，更是团队技术认知与协作
能力升级的“催化剂”。每一次难关的突破，都伴随开发规范的完善与技术思维的深化。
9.2.1 跨域与资源加载：从“被动规避”到“主动规范”
CORS 协议与浏览器资源加载限制是项目初期最棘手的问题。8 月 29 日，梁思齐
在调试地图矩阵（用数字标记交互区域）时，发现直接打开 HTML 文件无法读取 JS 配
43
项目计划
置，控制台报错“Access to script at ’file:///...’ from origin ’null’ has been blocked by CORS
policy”；同期，潘治鑫的存档系统也出现“localStorage 数据无法跨页面共享”的问题。
初期，我们尝试修改文件路径、改用绝对路径调用，均以失败告终。在 8 月 29 日晚间
会议中，我们转变思路：先调研往届解决方案（发现“瑞雪”项目用 Python HTTP 服务
器启动），再分工测试不同方案——潘治鑫尝试用 Python 的 http.server 模块（命令
“python -m http.server 8000”），成功解决跨域问题；梁思齐则测试 Live Server 插
件，发现能自动处理资源加载。但新问题随之而来：不同成员使用不同启动方式，导致
代码路径调用混乱。为彻底解决此问题，我们共同制定 “资源调用三规范”：一是所有
文件统一使用相对路径；二是建立“assets”目录分类管理 JS、CSS、图片资源；三是在
readme 文档中明确“Python HTTP 服务器与 Live Server 双启动方式”的操作步骤。此规
范不仅解决了当前问题，更避免了后续多模块协作时的路径冲突，体现 “问题解决 - 规
范沉淀”的技术思维升级。
9.2.2 多系统联动：从“各自为战”到“接口协同”
随着模块增多，“系统间数据同步”成为新挑战。9 月 11 日，梁思齐在对接战斗系
统与地图主系统时，发现战斗结果无法回传至主游戏，导致玩家完成战斗后卡在结算界
面。经排查，问题根源在于：战斗系统的结果参数仅在本地存储，未通过存档接口同步
至主游戏的状态；同时，主游戏未监听战斗页面的关闭事件，无法主动拉取最新数据。
针对此问题，我们成立专项攻坚小组（梁思齐、潘治鑫、房金衡），通过三次短时会议
确定解决方案：一是在战斗系统的“战斗结束”函数中，新增调用存档接口的逻辑；二
是主游戏页面添加“onmessage”监听，接收战斗页面的关闭通知后，拉取最新数据；三
是定义统一的 URL 参数格式，确保参数传递无歧义。此次攻关让我们深刻认识到 “接
口设计先行”的重要性。此后，我们在新增模块前，都会先在每日会议中确定接口参数，
避免后期返工。这种 “先定义接口、再开发功能”的思路，使后续修机小游戏与地图门
解锁的对接仅用 2 小时完成，效率提升 60%。
9.2.3 像素美术与技术适配：从“风格割裂”到“视觉统一”
美术与技术的适配问题贯穿开发全程。王宇豪作为美工，初期面临“素材风格不统
一”“像素比例失调”两大难题。为解决这些问题，我们建立 “美术 - 技术协同机制”：
一是素材标准统一，王宇豪明确“所有素材采用 32×32 像素、统一暗色调”，并用 GIMP
批量处理；二是交互区域可视化，梁思齐在 Canvas 绘制时，为可交互物件添加半透明红
色边框（测试模式下显示），方便王宇豪校准美术位置；三是效果迭代反馈，王宇豪为
实验笔记添加发光动画后，梁思齐同步调整触发范围，确保玩家能清晰感知交互点。最
终，地下室地图的视觉与技术适配效果显著：场景风格统一，交互准确率达 100%。此
次协作也让我们明白：美术不仅是“视觉装饰”，更是 “交互引导”——王宇豪后期为
关键道具添加的发光动画，使玩家找到实验笔记的平均时间从 2 分钟缩短至 30 秒，极
大提升游戏体验。
44
项目计划
9.3 协作认知升级：从“分工”到“共生”的团队蜕变
三周开发不仅是项目落地的过程，更是团队协作认知从“简单分工”向“深度共生”
跃迁的过程。我们深刻体会到：优秀的团队不是“各做各的”，而是“彼此支撑、互为补
充”。
9.3.1 每日会议：从“进度同步”到“问题共治”
我们坚持的“每日 19 点线上会议”，是团队协作的 “核心枢纽”。会议初期以“各
自说进度”为主，效率较低；后改革为 “问题优先”模式——每次会议先聚焦 3 个核心
问题，再同步进度、分配任务。这一改革带来显著变化：针对 CORS 跨域问题，会议仅
用 40 分钟就确定解决方案；讨论战斗难度梯度时，技术、策划、美术共同形成“章节 -
角色 - 数值 - 视觉”四维联动的难度设计。改革后的会议解决问题效率提升 150%。更重
要的是，每日会议培养了 “责任共担”的意识。这种“补位不缺位”的协作，让项目在
成员有临时任务时仍能高效推进。
9.3.2 分工与补位：从“固定角色”到“动态支撑”
项目初期，我们按“技术、策划、美术”划分固定角色，但发现“绝对分工”无法
应对复杂需求。为此，我们转向 “核心角色 + 动态补位”的模式：每个成员有明确核心
职责，但同时需掌握其他领域的基础技能，以便在团队需要时补位。这种模式成效显著：
潘治鑫（核心存档）协助优化移动逻辑，用“状态机”思路重构代码；王宇豪（核心美术）
学习基础 HTML/CSS，协助优化页面排版，使响应式适配率从 70% 提升至 100%。“动
态补位”不仅提升了项目效率，更让每个成员看到 “自身工作的价值关联”，形成良性
协作循环。
9.3.3 冲突与共识：从“意见分歧”到“最优解共建”
开发中并非毫无冲突，而是通过“理性讨论、数据支撑”将冲突转化为“共识升级”
的契机。关于“战斗系统是否需要难度梯度”的分歧，我们采取 “数据验证 + 最小原
型”的方案：房金衡计算战斗回合数，梁思齐快速制作“配置表原型”，最终团队达成共
识——采用“配置表 + 动态加载”方案，既实现难度梯度，又避免代码冗余。此次冲突
让我们建立 “冲突解决三原则”：一是不回避分歧；二是用数据说话；三是做最小原型。
这一原则确保团队决策始终基于“可行性”与“玩家体验”，而非个人偏好。
9.4 项目开发认知：从“功能堆砌”到“用户体验”的思维转变
三周开发让我们对“游戏开发”的认知实现从“技术驱动”到“用户为中心”的深
刻转变。我们逐渐明白：优秀的游戏不仅是“功能的集合”，更是 “体验的闭环”。
9.4.1 从“实现功能”到“优化体验”：细节决定成败
初期，我们更关注“功能是否实现”，而忽略“体验是否流畅”。例如，角色移动在高
刷新率电脑上过快，成就弹窗会遮挡战斗界面。针对测试反馈，我们启动 “体验优化专
项”：梁思齐用“requestAnimationFrame”替代固定时间间隔，并添加移动加速度，使
移动手感更佳；刘雨桐优化成就弹窗，新增“战斗中延迟显示”逻辑，并调整位置至右
45
项目计划
下角。体验优化后的效果显著：第二次测试中，玩家对“移动流畅度”的满意度从 40%
提升至 90%。这让我们深刻认识到：“能用”与“好用”的差距在细节。
9.4.2 从“开发者视角”到“玩家视角”：换位思考的重要性
开发初期，我们常陷入“开发者视角”的误区，默认玩家“懂我们的设计逻辑”。例
如，地图初期无任何引导，导致 70% 的测试者找不到门。为解决此问题，我们建立“玩
家视角验证机制”：一是引入外部测试者；二是推行“自我盲测”；三是简化信息传递
（添加“剧情回顾”按钮和“光效引导”）。“玩家视角”的转变带来显著效果：外部测试
者找到实验室的平均时间从 15 分钟缩短至 3 分钟。这一过程让我们明白：游戏开发不
是“自我表达”，而是 “与玩家的沟通”。
9.4.3 从“单打独斗”到“工程化思维”：项目管理的价值
项目初期，我们缺乏工程化思维，导致代码无规范、文档无记录、版本无管理。为
解决这些问题，我们逐步建立 “工程化开发体系”：一是制定《前端开发规范》，统一代
码风格；二是推行文档管理，每个模块完成后编写技术文档；三是采用“日期 + 版本号”
的命名规则进行版本控制。工程化体系的建立，使项目后期维护效率大幅提升。这让我
们认识到：小型项目同样需要工程化思维——规范不是“束缚”，而是“效率保障”；文
档不是“多余工作”，而是“团队协作的桥梁”。
9.5 总结与展望：不止于“完成作业”，更是“能力成长”
三周《第五人格：推演终幕》的开发，对我们而言不仅是完成课程作业，更是一次
完整的“小型项目实战”。我们收获的不仅是一个可运行的游戏作品，更是技术能力、协
作意识、项目思维的全方位成长。
回顾整个开发过程，我们也清醒地认识到不足：
• 原创美术能力不足：王宇豪虽完成基础素材，但原创角色立绘仍有提升空间。
• 游戏可玩性深度不够：战斗系统虽实现基础卡牌玩法，但缺乏“卡组流派”“角色
组合策略”等深度设计。
• 工程化体系仍需完善：未使用 Git 进行版本控制，多成员协作时仍有冲突风险。
未来，若有机会迭代项目，我们计划从三方面改进：
• 提升美术原创性：学习 Aseprite 高级绘画技巧，为每个角色设计独特的战斗动作。
• 深化玩法设计：新增“卡组预设”“角色羁绊”系统，提升战斗策略性。
• 完善工程化工具：引入 Git 进行版本控制，使用 Figma 进行 UI 设计协作，进一步
提升团队开发效率。
此次项目开发，让我们深刻体会到：“做项目”是最好的学习方式。这三周的经历，
将成为我们后续学习与工作中宝贵的财富，激励我们继续在“技术与协作”的道路上不
断探索、不断成长。